% =====================================================================
%  1bat-ud4-10.tex  —  SESSIÓ 10: Radiografia d'un barri (piràmides de població)
%  (sense preàmbul: s'inclou des de main.tex)
% =====================================================================
\horatitol{Sessió 10 — Radiografia d'un barri (piràmides de població)}%
{Llegir dues piràmides reals, comparar un territori envellit amb un de jove i redactar un diagnòstic d'equipaments.}

\subsection{Idea clau}

A la sessió 9 vam aprendre a llegir gràfics amb mirada crítica. Ara ho apliquem al
gràfic més útil per a un tècnic de serveis a la comunitat: la \textbf{piràmide de
població}. Llegint-la, podem fer una \emph{radiografia} d'un territori i deduir
\textbf{quins serveis socials necessitarà}.

\begin{idea}
\textbf{Com es llegeix una piràmide de població}
\begin{itemize}[leftmargin=1.4em, itemsep=2pt]
  \item \textbf{Eix vertical:} les \textbf{franges d'edat} (les més grans a dalt).
  \item \textbf{Eix horitzontal:} el \textbf{percentatge} de població de cada franja;
        \textbf{homes} a l'esquerra, \textbf{dones} a la dreta.
  \item \textbf{Base ampla} $\rightarrow$ població \textbf{jove} (molts infants).
  \item \textbf{Cim ample / base estreta} $\rightarrow$ població \textbf{envellida}.
  \item La \textbf{moda} és la franja d'edat \emph{més nombrosa} (la barra més llarga).
\end{itemize}
\textbf{De la forma al servei:} la forma de la piràmide anticipa quins equipaments
caldran.
\end{idea}

\subsection{Exemples}

\begin{exemple}[un municipi envellit]
Aquesta piràmide correspon a un municipi rural amb població molt envellida. Fixa't que
la \textbf{base és estreta} (pocs infants) i el \textbf{cim és ample} (molta gent
gran):

\begin{center}
\begin{tikzpicture}
\begin{axis}[
    width=11cm, height=6.4cm,
    xbar, bar width=8pt, y=14pt,
    xmin=-16, xmax=16,
    xtick={-16,-12,-8,-4,0,4,8,12,16},
    xticklabels={16,12,8,4,0,4,8,12,16},
    xticklabel style={font=\scriptsize, color=udgrey},
    xlabel={\small \% població \quad (homes $\leftarrow$ \ $\rightarrow$ dones)},
    ytick={1,2,3,4,5,6},
    yticklabels={0-14,15-29,30-44,45-64,65-79,80+},
    yticklabel style={font=\scriptsize, color=udgrey},
    tick style={udgrey}, axis line style={udgrey},
    enlarge y limits=0.12,
    legend style={font=\scriptsize, at={(0.5,-0.28)}, anchor=north, legend columns=2}]
  \addplot[fill=udblue!55, draw=udnavy] coordinates {(-4,1) (-5,2) (-6,3) (-10,4) (-13,5) (-8,6)};
  \addplot[fill=udblue!25, draw=udnavy] coordinates {(4,1) (5,2) (6,3) (11,4) (14,5) (12,6)};
  \legend{homes, dones}
\end{axis}
\end{tikzpicture}

\figcap{Municipi envellit: la franja més nombrosa (moda) és la de $65$–$79$ anys; molt poca base.}
\end{center}

\noindent La franja \textbf{més llarga} és la de $65$–$79$ anys: és la \textbf{moda}.
Hi ha molta gent gran i pocs infants: un municipi \textbf{envellit}.
\end{exemple}

\begin{exemple}[un barri jove amb molta immigració]
Aquesta segona piràmide és d'un barri amb població jove. La \textbf{base és ampla}
(molts infants i joves) i el \textbf{cim, estret}:

\begin{center}
\begin{tikzpicture}
\begin{axis}[
    width=11cm, height=6.4cm,
    xbar, bar width=8pt, y=14pt,
    xmin=-18, xmax=18,
    xtick={-16,-12,-8,-4,0,4,8,12,16},
    xticklabels={16,12,8,4,0,4,8,12,16},
    xticklabel style={font=\scriptsize, color=udgrey},
    xlabel={\small \% població \quad (homes $\leftarrow$ \ $\rightarrow$ dones)},
    ytick={1,2,3,4,5,6},
    yticklabels={0-14,15-29,30-44,45-64,65-79,80+},
    yticklabel style={font=\scriptsize, color=udgrey},
    tick style={udgrey}, axis line style={udgrey},
    enlarge y limits=0.12,
    legend style={font=\scriptsize, at={(0.5,-0.28)}, anchor=north, legend columns=2}]
  \addplot[fill=udblue!55, draw=udnavy] coordinates {(-15,1) (-16,2) (-14,3) (-9,4) (-4,5) (-2,6)};
  \addplot[fill=udblue!25, draw=udnavy] coordinates {(14,1) (15,2) (13,3) (9,4) (5,5) (2,6)};
  \legend{homes, dones}
\end{axis}
\end{tikzpicture}

\figcap{Barri jove: la franja més nombrosa (moda) és la de $15$–$29$ anys; base molt ampla.}
\end{center}

\noindent Aquí la franja \textbf{més llarga} és la de $15$–$29$ anys, i la base
($0$–$14$) també és molt ampla: un barri \textbf{jove}.
\end{exemple}

\subsection{Exercicis resolts}

\begin{exresolt}[comparar les dues formes]
Compara la piràmide del municipi envellit amb la del barri jove: identifica la moda de
cadascuna i digues què anticipa cada forma sobre la població futura.
\solucio
\textbf{Municipi envellit:} moda a $65$–$79$ anys; base estreta. Forma de
\emph{piràmide invertida} (ampla a dalt). Anticipa una població que \textbf{continuarà
envellint}: cada cop hi haurà més gent gran i menys jovent que la substitueixi.

\textbf{Barri jove:} moda a $15$–$29$ anys; base ampla. Forma de \emph{piràmide
clàssica} (ampla a baix). Anticipa una població \textbf{en creixement}, amb molts
infants que en pocs anys seran adolescents i joves.
\resposta{Envellit: moda $65$–$79$, seguirà envellint. Jove: moda $15$–$29$, base
ampla, població en creixement.}
\end{exresolt}

\begin{exresolt}[diagnòstic d'equipaments]
A partir de les dues piràmides, redacta un breu diagnòstic: quins equipaments i
serveis socials caldran a cada territori els pròxims deu anys?
\solucio
\textbf{Municipi envellit} (molta gent gran): caldran \textbf{casals d'avis},
\textbf{atenció domiciliària}, serveis de \textbf{teleassistència}, places de
\textbf{centre de dia} i \textbf{residència}, i transport adaptat. La prioritat és
l'atenció a la dependència.

\textbf{Barri jove} (molts infants i joves): caldran \textbf{llars d'infants} i
\textbf{escoles}, \textbf{casals i activitats juvenils}, programes d'\textbf{ocupació
juvenil} i, atesa la immigració, serveis d'\textbf{acollida i integració}. La
prioritat és l'educació i la inserció.
\resposta{Envellit $\rightarrow$ serveis de dependència (casals d'avis, atenció
domiciliària). Jove $\rightarrow$ educació i ocupació (llars d'infants, ocupació
juvenil, acollida).}
\end{exresolt}

\subsection{Exercicis per fer}

\begin{exercicis}
\begin{exlist}
  \item Mira la piràmide del \textbf{municipi envellit}. Quina és la franja d'edat amb
        més població (la moda)? És més ampla la base o el cim?
  \item Mira la piràmide del \textbf{barri jove}. Quina és la moda? Quina forma té la
        base?
  \item En la piràmide del municipi envellit, compara la franja $0$–$14$ amb la
        $65$–$79$. Quina té més població? Què vol dir això sobre el relleu
        generacional?
  \item \textbf{(Diagnòstic)} Per al \textbf{barri jove}, escriu \textbf{tres}
        equipaments o serveis socials que prioritzaries i justifica per què, mirant la
        piràmide.
  \item \textbf{(Interpretació)} Un poble té una piràmide gairebé rectangular (totes
        les franges semblants) excepte la de $80+$, que és molt ampla. Com el
        descriuries? Quin servei hi seria especialment necessari?
  \item \textbf{(Raonament)} Per què és tan útil, per a la planificació de serveis
        socials, mirar la piràmide \textbf{abans} de decidir on s'inverteix? Posa un
        exemple d'un error que es podria cometre si no es mira.
\end{exlist}
\end{exercicis}

% ---------------------------------------------------------------------
\fullsolucions{Sessió 10}
\begin{solucions}
\begin{sollist}
  \item La moda és la franja \textbf{$65$–$79$ anys} (la barra més llarga). El
        \textbf{cim} és més ample que la base (típic d'una població envellida).
  \item La moda és la franja \textbf{$15$–$29$ anys}. La base ($0$–$14$) és
        \textbf{molt ampla} (molts infants).
  \item La franja $65$–$79$ té \textbf{molta més població} que la $0$–$14$. Vol dir que
        \textbf{no hi ha relleu generacional}: hi ha molta més gent gran que infants
        que els substitueixin, de manera que el municipi continuarà envellint.
  \item Resposta oberta. Exemples vàlids: llars d'infants (per la base ampla de
        $0$–$14$), programes d'ocupació juvenil (per la moda a $15$–$29$), serveis
        d'acollida i integració (per la immigració), casals juvenils i suport educatiu.
  \item Seria una població \textbf{adulta i molt envellida alhora}: poca natalitat però
        un pes enorme de la gent molt gran ($80+$). Hi serien especialment necessaris
        els serveis d'\textbf{atenció a la dependència} (atenció domiciliària,
        residència, teleassistència).
  \item Resposta oberta. Idea: la piràmide diu quines necessitats té realment el
        territori; invertir sense mirar-la pot portar a, p.\,ex., construir més llars
        d'infants en un municipi envellit (on gairebé no hi ha nadons) en comptes de
        places d'atenció a la gent gran, que és el que de debò es necessita.
\end{sollist}
\end{solucions}
